% [메모] 8/30 v2 개정 (가이드라인 [Conclusion] 3단락 + WRITING_BRIEF §4,8,9).
%   1단락 문제·제안, 2단락 실측 요약(시뮬 3물체 — laptop/실물은 \rev),
%   3단락 limitation: 채택 요소 3개 + 실험이 드러낸 2개(부피 오차의 조건수
%   비용, 관성 텐서 비식별) + future work(관절각 dynamics, RPNA).
\section{Conclusion}

A sim-ready articulated asset requires per-part volume \emph{and} density;
vision-based estimation is inaccurate, and interaction-based identification
rigid-only. PIVoT identifies per-part density in a closed loop:
information-gain joint-configuration selection, three gravity-direction
measurement poses linked by collision-free paths, and uncertainty-aware
updates.
% [번역] 관절 물체의 sim-ready asset 에는 part 별 부피'와' 밀도가 모두
%   필요한데, 비전 기반 추정은 부정확하고 상호작용 기반 식별은 강체에
%   국한되어 왔다. 우리는 정보이득 기반 관절 형상 선택, 서로 다른 중력 방향의
%   측정 자세 3개와 그 사이 collision-free 경로, 불확실성 기반 밀도 갱신으로
%   구성된, part 별 밀도를 능동 상호작용으로 식별하는 폐루프 파이프라인
%   PIVoT 를 제시했다.

On three simulated objects, worst-part density error was 0.18, 0.11, and
0.23\%, versus a homogeneous baseline's 49.1, 321.0, and 34.6\%. Stronger
evidence is the interval: our 95\% intervals covered the truth in 62 of 64
part--seed pairs, a WLS ablation in only 33 of 64---the errors-in-variables
bias does not shrink with rounds, and a WLS interval cannot see it---and a
single configuration cannot certify five unknowns, its half-widths reaching
42.5\% (Cor.~\ref{cor:rigid}). Omitting the custom objects' friction hinge
as its own unknown yields a 79.1\% hinge error under a claimed 0.83\%
half-width. The intervention added for physics estimation is one
joint-configuration adjustment; the reconstruction-stage intervention is
inherited from \cite{lee2026rora}. \rev{Validation on the laptop asset and
the physical robot is in preparation.}
% [번역] 시뮬레이션 3개 물체에서 최악 part 의 상대 밀도 오차는 0.18, 0.11,
%   0.23\% 로, 균일밀도 baseline 의 49.1, 321.0, 34.6\% 와 대비된다. 더 강한
%   근거는 구간이다. 우리의 95\% 구간은 64개 부위-시드 쌍 중 62를 덮었고
%   WLS ablation 은 33에 그쳤다 — errors-in-variables 치우침은 라운드를
%   늘려도 줄지 않고, WLS 구간은 그것을 보지 못한다 — 단일 형상은 미지수
%   다섯을 증명하지 못해 반폭이 42.5\% 에 이른다(따름정리). 마찰 힌지를 독립
%   미지수로 두지 않으면 반폭 0.83\% 를 주장하면서 힌지 오차 79.1\% 를 낸다.
%   물성 추정을 위해 추가되는 개입은 관절 형상 조정 한 번이며, 재구성 단계의
%   개입은 [RORA] 에서 상속된다. (laptop 자산과 실물 로봇 검증은 준비 중 —
%   \rev{} 표시.)

Limitations begin with adopted components: the input reconstruction's
segmentation and joint errors propagate downstream, with its
human-in-the-loop intervention \cite{lee2026rora}; sampling-based RRT grows
costly on complex topologies without optimality guarantees; and a single
linkage gripper on a 6-DoF arm restricts validation to small, light
objects. The experiments add two. Volume error is not benign: density error tracks
it one-to-one, inertia as $s^2$ (Rem.~\ref{rem:volume}), part mass is far
less sensitive yet not exactly invariant, and it multiplies rounds
(Sec.~\ref{sec:vol})---the practical limit is conditioning, set by volume
and joint-angle precision, not rank. The inertia tensor is not identified: the quasi-static
wrench never sees it, so we derive it from mesh geometry and
identified density, inheriting mesh scale error quadratically. Future work: estimate joint-angle dynamics and reorient accordingly, where
tools like RPNA~\cite{wensing2024geometric} become
discriminative.
% [번역] 한계는 먼저 채택한 구성 요소에서 온다. 입력 재구성의 part
%   segmentation 과 관절 추정 오류는 human-in-the-loop 개입과 함께 이후
%   단계로 상속되고 [RORA], 샘플링 기반 RRT 는 복잡한 topology 에서 계산이
%   늘며 최적성을 보장하지 않고, 단일 linkage gripper 와 6자유도 팔은 검증
%   대상을 소형·경량 물체로 좁힌다. 실험은 둘을 더한다. 부피 오차는 무해하지
%   않다. 밀도 오차는 이를 1:1 로, 관성은 s^2 로 따라가고(비고), part 질량은
%   훨씬 덜 민감하지만 엄밀히 불변은 아니며, 라운드 수를 곱으로 늘린다
%   (IV-B: 1→7, 3/8 미수렴). 즉 실질적 한계는 rank 가 아니라 부피·관절각
%   정밀도가 정하는 조건수다. 또한 관성 텐서는 식별되지 않는다. 준정적 회전 방정식에서
%   사라지므로 메시 기하와 식별된 밀도에서 유도하며, 메시 배율 오차를 제곱으로
%   물려받는다. Future work 는 관절각 dynamics 를 추정해 그에 맞춰 물체를
%   재배향하는 것이며, 거기서는 RPNA 같은 도구가 실제 판정력을 갖는다.
